% novices.tex                
% Note, that this a HyperLaTeX source and not plain LaTeX!
\documentclass[a4paper,11pt,twoside,titlepage,dvips]{scrartcl}
\usepackage{hyperlatex}
\usepackage{a4wide}
\usepackage{times}
\usepackage[latin1]{inputenc}
\usepackage[T1]{fontenc}
\usepackage{graphicx}
\usepackage{alltt}
\usepackage{moreverb}
\usepackage{fancyhdr}
\W\usepackage{rhxpanel}
\W\usepackage{sequential}


\T\DeclareGraphicsExtensions{.eps.gz,.eps}


% Hyperref should be among the last packages loaded
\usepackage{hyperref}

% Macros specific to this package
\newcommand{\Button}[1]{[\,#1\,]}
\newcommand{\Menu}[1]{\emph{#1}}
\newcommand{\Filename}[1]{\texttt{#1}}

% Some people still insist that Email should be written as two words
% for whatever reasons.  Maybe is is that they don't like the game of
% Teekesselchen and fear the ambiguity between email and enamel (which
% uses the same word in German).  

\newcommand{\Email}{email}

\newcommand{\EchelonUrl}{http://www.heise.de/tp/r4/artikel/6/6928/1.html}
\newcommand{\EinsteigerPDFURL}{http://wald.intevation.org/frs/download.php/134/gpg4win-fuer-einsteiger-2.0.0.RC1.pdf}
\newcommand{\DurchblickerPDFURL}{http://wald.intevation.org/frs/download.php/134/gpg4win-fuer-durchblicker-2.0.0.RC1.pdf}

\newcommand{\OtherBook}{\texorhtml{$\spadesuit$}{\xml{p}==\xmlent{gt}}}

\newcommand{\IncludeImage}[2][]{\texorhtml{%
\includegraphics[#1]{#2}%
}{%
\htmlimg{#2.png}%
}}

\newcommand{\HlxIcons}{}

%%\W\htmlpanelgerman
%%\W\extrasgerman
%%\W\dategerman
%%\W\captionsgerman
%%\htmlpanelfield{Inhalt}{hlxcontents}
\W\renewenvironment{thefootnotes}{\chapter{Footnotes}
\W  \begin{description}}{\end{description}}



\newcommand{\manualversion}{\manualversionNovices}
\newcommand{\manualdate}{\manualdateNovices}

\T\fancyhead{} % clear all fields
\T\fancyhead[LO,RE]{Gpg4win for Novices \manualversion\ \manualinprogressEn}
\T\fancyhead[RO,LE]{\thepage}
\T\fancyfoot[C]{\includegraphics[width=1cm]{gpg4win-logo}}


% Title stuff 
\htmltitle{Novices}
%\htmladdress{Gpg4win Project, \today}
\title{
\IncludeImage[width=8cm]{gpg4win-logo}
\\
Gpg4win for Novices}

\author{\htmlonly{\xml{p}\small
Hier \xlink{Download als PDF Version 2.0.1 (1.5
MByte)}{http://wald.intevation.org/frs/download.php/203/gpg4win-fuer-einsteiger-2.0.1.pdf} und
\xlink{Download�bersicht alle PDF
Versionen}{http://wald.intevation.org/frs/?group_id=11}\xml{br}
Zu \xlink{Gpg4win for Advanced Users}{advanced.html}\xml{br}
Zur \xlink{Gpg4win Homepage}{http://www.gpg4win.de/}\xml{p}
}%
A publication of the Gpg4win project\\
  \small based on original documents by  
\T\\
  \small Manfred J. Heinze, Karl Bihlmeier, Isabel Kramer
\T\\
  \small Dr. Francis Wray und Ute Bahn
\T\\ \
  \small Werner Koch
 \\
  \small Translated from the German original by 
\T\\
  \small Brigitte Hamilton}
\date{Version \manualversion\ as of \manualdate\ \manualinprogressEn}


\begin{document}
\thispagestyle{empty}
\pagestyle{fancy}
\T\parindent0cm
\T\parskip\medskipamount


\maketitle


\section*{Impressum Gpg4win}

\noindent
Copyright \copyright{} 2002 Bundesministerium f�r Wirtschaft und
Technologie\\
Copyright \copyright{} 2005 g10 Code GmbH\\
Copyright \copyright{} 2006 Brigitte Hamilton\\
Permission is hereby granted to copy, distribute and/or modify this
document
under the terms of the GNU Free Documentation License, Version 1.2 (or
newer)
published by the Free Software Foundation; with
the Invariant Sections being "`Impressum"', no Front-Cover Texts, and
no Back-Cover Texts.  A copy of the license is included in the section
entitled "`GNU Free Documentation License"'.

%%\htmlonly{Die aktuelle PDF Version dieses Dokuments finden sie unter
%%\xlink{\EinsteigerPDFURL}{\EinsteigerPDFURL}.}

\vfill

\begin{center}
  \T\fbox{\parbox{.7\textwidth}{%
      The statements on the \textbf{next page} are not anymore
      correct.  Due to a wrong application of the terms of
      the GFDL, it is legally not possible to fix them.  Please add
      new copyright notices only here.  \T}}
\end{center}

\clearpage
%% Orginal page 4
%% This page 4 is an invariant section. Translation not allowed.
\section*{Impressum}
Diese Seite darf nicht ver�ndert werden.\\

Autor: Manfred J. Heinze, TextLab text+media\\
Beratung: Lutz Zolondz, G-N-U GmbH\\
Illustrationen: Karl Bihlmeier, Bihlmeier \& Kramer GbR\\
Layout: Isabel Kramer, Bihlmeier \& Kramer GbR\\
Fachtext: Dr. Francis Wray, e-mediate Ltd.\\
Redaktion: Ute Bahn, TextLab text+media\\
1. Auflage, M�rz 2002\\


Copyright \copyright{} Bundesministerium f�r Wirtschaft und Technologie

Dieses Buch unterliegt der "`GNU Free Documentation License"'.
Originaltext der Lizenz: http://www.\texonly{\linebreak{}}gnu.org/copyleft/fdl.html.
Deutsche �bersetzung http://nautix.sourceforge.net""/docs/fdl.de.html
sowie auf der beiliegenden CD-ROM.  Es wird die Erlaubnis gegeben,
dieses Dokument zu kopieren, zu verteilen und/oder zu ver�ndern unter
den Bedingungen der GNU Free Documentation License, Version 1.1 oder
einer sp�teren, von der Free Software Foundation ver�ffentlichten
Version.  Diese Seite ("`Impressum"') darf nicht ver�ndert werden und
muss in allen Kopien und Bearbeitungen erhalten bleiben
("`unver�nderlicher Abschnitt"' im Sinne der GNU Free Documentation
License).  Wenn dieses Dokument von Dritten kopiert, verteilt und/oder
ver�ndert wird, darf in keiner Form der Eindruck eines Zusammenhanges
mit dem Bundesministerium f�r Wirtschaft und Technologie erweckt
werden.  Wie das OpenSource-Kryptografieprogramm GnuPP selbst wurden
diese Texte nicht f�r Mathematiker, Geheimdienstler und Kryptografen
geschrieben, sondern f�r jedermann.



\clearpage %% End of original page 4.

\tableofcontents

%%\clearpage
%% Orginal page  6
%% We don't use these foreword anymore because Mr. M�ller is not
%% anymore minister of economic and technology.  We might want to ask
%% for a new foreword by the current head of that ministr
%%

\clearpage
%% Orginal page 7
\section{About this manual}

The Gpg4win manual and exercise module consists of three parts:

\begin{itemize}

\item \textbf{the condensed version "`Gpg4win for Novices"',} which you
are reading right now,

\item \textbf{the manual  "`Gpg4win for Advanced Users"'} in PDF-Format,
which can be found on your hard drive following the installation of
Gpg4win,

\item \textbf{the exercise ROBOT Adele,} which allows you to practice
the \Email{} en- and decryption process(an internet connection is
required). 
\end{itemize}


\textbf{"`Gpg4win for Novices"'} provides you with a quick guide to the
installation and everyday use of the Gpg4win software. The time required
to work thru this manual is about half an hour, but will depend on your
knowledge of computers and Windows. 

\textbf{"`Gpg4win for Advanced Users"'} provides more in-depth
information on the basic principles and mechanisms behind Gpg4win as
well as its less commonly used capabilities.

Both manuals are available in PDF format, so you can print your own
manual if you did not receive a printed copy. 

Each manual can be used independently, however, we suggest that you read
both manuals to get a better understanding of the software.

\OtherBook{} This symbol references a link to the other manual.


\textbf{The practice robot Adele} is available on the Internet. Adele
receives, sends an decrypts encrypted \Email{}s, and is a useful tool
which will assist you in becoming more familiar with the software. 

Adele was developed as part of the older GnuPP project, where it is
still used.  "`Gpg4win for Novices"' also uses this very reliable
practice robot and hereby wishes to express its gratitude to the owners
of gnupp.de for the operation of Adele. 

\clearpage
%% Orginal page 8
\section{What is Gpg4win?}

% what-is-gpg4win.tex  - Introduction used by both manuals


\textbf{
Gpg4win (GNU Privacy Guard for Windows) is an email encryption software.
It is the result of a project initiated by the Federal Office for
Information Security, and includes the following components: 
}

\textbf{GnuPG}: its key component, the encryption software \newline
\textbf{GPA}: GNU Privacy Assistant,
a key manager \newline
\textbf{WinPT}: Key Manager, which also supports encryption via your
Clipboard \newline
\textbf{GPGol}: a plugin for Microsoft Outlook which integrates the
operation of GnuPG \newline
\textbf{GPGee}: a plugin for Windows Explorer which allows encryption of
data by right-clicking on your mouse \newline
\textbf{Sylpheed-Claws}: a complete email program with integrated GnuPG
operation \newline

The encryption program GnuPG (GNU Privacy Guard) provides you with a
secure, simple and free method of email encryption. It can be used
privately or commercially without restrictions.  The encryption
technology used by GnuPG is extremely secure and cannot be broken using
current technology.

GnuPG is a free software\footnote{sometimes incorrectly identified as
Open Source Software}. This means that anyone can use the software for
private or commercial purposes, as well as analyze or change the source
codes (ie. the actual programming commands), and distribute the
same.\footnote{However, keep in mind that a fair amount of technical
knowledge is required to change the program, otherwise the program's
security may be compromised.}

The transparency of the source code forms an essential part of a
security software, as it is the only way to verify the trustworthiness
of the program.


GnuPG is based on the international standard OpenPGP (RFC 2440), is
fully compatable with PGP and uses the same infrastructure (key server
etc.).

PGP ("`Pretty Good Privacy"') is not free software; many years ago it
was available on a temporary basis under similar conditions as GnuPG,
but is no longer considered state-of-the-art.  

Additional information regarding GnuPG and other projects undertaken by
the german government in the area of Internet security can be found on
the website of the Federal Office for Information Security \W
\xlink{www.bsi-fuer-buerger.de}{http://www.bsi-fuer-buerger.de}
\T\href{http://www.bsi-fuer-buerger.de}{www.bsi-fuer-buerger.de}





\clearpage
%% Orginal page 9
\section{Installing Gpg4win}

If you already have a GnuPG-based application such as GnuPP, GnuPT,
WinPT or GnuPG Basics installed on your system, we recommend reading the
Appendix \ref{ch:migration} regarding the transferability of existing
keys. 

Installing Gpg4win from a CD-ROM:

Insert the CD-ROM in the CD-ROM drive of your computer and log in as
Administrator. On your screen click on the CD-ROM icon titled 'Gpg4win'.
Once the CD-ROM icon opens, click on the installation icon titled
'Gpg4win'.

Installing Gpg4win from the Internet:

If you have downloaded Gpg4win from the Internet, click on the new file
(it should be named \texttt{gpg4win-\PackageVersion{}.exe}, or a newer
version).  Please ensure that you have downloaded the file from a
trustworthy site.

The following installation steps apply in all situations:

You will be asked if you want to install the program; click on
\Button{Yes}. 

The following screen will appear:

% screenshot:  Welcome Seite Installer
\begin{center}
\IncludeImage{sc-inst-welcome}
\end{center}

If you have other programs running on your computer, close them now and
click on \Button{Next}.

\clearpage
%% Orginal page 10

The licencing page contains information regarding the licencing of this
software.

If your sole intention is to install and use the software, you are not
obliged to read this information. 

However, if you plan to distribute and/or alter the software, you must
familiarize yourself with the conditions contained in the licencing
agreement.


% Screenshot Lizenzseite des Installers
\begin{center}
\IncludeImage{sc-inst-license}
\end{center}


Click on \Button{Next}.


\clearpage
%% New page (not in original document)

On the components page you can select which features of GnuPG For
Windows you want to install. 
 
To assist you with your selection, a short description dialogue appears
when you roll your mouse over the selected item. At this point, you may
also want to check the available memory on your hard drive.

% sreenshot Auswahl zu installierender Komponenten
\begin{center}
\IncludeImage{sc-inst-components}
\end{center}

The recommended minimum installation consists of GnuPG, GPA, WinPT and
the manuals. The remaining programs can be installed later as required. 

Click on \Button{Next}.


\clearpage
%% Original page 11

On this page you can choose the folder in which to install Gpg4win on
your computer. If you do not specify a folder, the default location for
this folder will be as shown: \\
\Filename{C:$\backslash$Programme$\backslash$GNU$\backslash$GnuPG}\\

% screenshot: Auswahl des Installationsverzeichnis.
\begin{center}
\IncludeImage{sc-inst-directory}
\end{center}

Click on \Button{Next}.

\clearpage

This page allows you to set start links for the program. The default
setting links Gpg4win to the start menu only.  Please note that these
settings can also be changed anytime within Windows.

% screenshot: Auswahl des Links
\begin{center}
\IncludeImage{sc-inst-options}
\end{center}

Click on  \Button{Next}.

\clearpage
%% Original page 12

If you chose to link the program with the Start Menu (as per previous
page), this page allows you to choose a Start Menu folder for the
program.


% screenshot:  Startmenu ausw�hlen
\begin{center}
\IncludeImage{sc-inst-startmenu}
\end{center}

For a standard installation select the default setting and click on
\Button{Install}.

\clearpage

During the installation, a progress bar will be displayed, as well as an
indicator listing the file currently being installed.  You can click on
\Button{Show details} to get more information on the installation steps
in progress.

% Screenshot: Ready page Installer
\begin{center}
\IncludeImage{sc-inst-ready}
\end{center}

Once the installation is complete, click on \Button{Next}.

\clearpage
%% Original page 13

The following page shows the last step of the installation process: 

% Screenshot: Finish page Installer
\begin{center}
\IncludeImage{sc-inst-finished} 
\end{center}

Click on  \Button{Finish}.

\clearpage

You may need to restart Windows for the settings to take effect.  In
that case, instead of showing the installation completion page, the
following page will appear: 

% Screenshot: Finish page Installer with reboot
\begin{center}
\IncludeImage{sc-inst-finished2}
\end{center}

At this point you can choose to restart Windows automatically or restart
later manually.

Click on \Button{Finish}.


% FIXME:  Wir m�ssen erkl�ren wie man Word als Standard Editor in
% Outlook ausschaltet.

\clearpage
%% Original page 14

\textbf{And that's it!}


You have now installed Gpg4win and are ready to work with the program.

Prior to starting Gpg4win, we recommend reading Chapter 3 and 4 of the
manual "`Gpg4win for Advanced Users"'
(PDF-Files).  These chapters highlight the ingeneous theory behind
Gpg4win's ability to encrypt \Email{}s in a safe and user-friendly
manner.

Of course you do not need to know the mechanics behind Gpg4win to be
able to use the program. However, because unlike other applications
Gpg4win is used to handle your most sensitive correspondence, it is a
good idea to understand the theory behind the program. 

$\ldots$

The following section will provide you with some tips for creating a
secure yet easy-to-remember passphrase.

\textbf{
\OtherBook{}  At this point, please read Chapter 3 and 4 in the manual
"`\xlink{Gpg4win for Advanced Users}{durchblicker.html}"'before reading
on.}

\clearpage
%% Original page 15
\section{Creating a key pair}

After reading the information explaining the security features of
Gpg4win as well as the creation of a good passwordphrase protecting your
private key, this section shows you how to create a key pair.

\textbf{
The process of creating a key, encryption and decryption is very
important, therefore it is a good idea to practice these steps until you
are very comfortable with the procedure. $\ldots$
}

The idea behind doing a "`dry-run"' is to help you gain confidence in
using the program, so that some of the more intricate steps in creating
keys will not pose any problems later on. To do this, you can use Adele,
a test server which was developed as part of a previous GnuPP project.
Adele allows you to try and test several practice keys before creating
your own set.  

\clearpage
%% Original page 16
\textbf{Let's get started!}

Start the GPA program from your Windows Start menu: 

% screenshot Startmenu with GPA highlighted
\begin{center}
\IncludeImage[width=0.4\textwidth]{sc-gpa-rungpa}
\end{center}

The following window will appear: 

% screenshot: GPA erste Schl�sselerzeugung
\begin{center}
\IncludeImage[width=0.9\textwidth]{sc-gpa-nokey}
\end{center}

Click on  \Button{Generate key now}.


\clearpage
%% Original page 17

\textbf{For practice purposes} you can enter any name in the window for
now,
i.e. "`Larry Smith"'.

% screenshot: Key generation wizard first page
\begin{center}
\IncludeImage[width=0.6\textwidth]{sc-gpa-gen-name}
\end{center}

Or you can generate your 'real' key in which case you would enter your
real name.

When you are done, click on  \Button{Forward}.

\clearpage
%% Original page 18
Now enter your \Email{} address.

Again, if this is a test run, you can use an imaginary \Email{}address
such as 
"`\verb-larrys@gpg4win.de-"'

% screenshot: Wizard email address
\begin{center}
\IncludeImage[width=0.6\textwidth]{sc-gpa-gen-email}
\end{center}

Or you can enter your real \Email{}address, and click on
\Button{Forward}.

%note by translator Brigitte: when I ran Gpg4win, the option outlined in
the text below was not available!

This option allows you to enter notes regarding your key.  Usually this
field is empty; however, if you are creating a test key, you should make
a note, such as "`test"'. The notes are a part of your User-ID and will
be made public along with your name and \Email{} address.  Now click on
\Button{Forward}.


\clearpage
%% Original page 19

Entering a passphrase is one of the most important steps in the program,
as the program is only as good (and secure) as your passphrase! 

Chapter 4 ("`The Passphrase"') from the manual "`Gpg4win for Advanced
Users"' already provided you with suggestions on creating a secure
passphrase.  So at this point you should have developed your own
passphrase - one that is private, easy to remember, and difficult to
crack.  

Enter your passphrase now. 

% screenshot: Wizard, Enter passphrase
\begin{center}
\IncludeImage[width=0.6\textwidth]{sc-gpa-gen-passwd}
\end{center}

If the passphrase you entered is not very secure, the program will give
you a warning in this regard.  This allows you to enter a more secure
passphrase. 

\textbf Again, you can enter a test passphrase for practice purposes, or
your real passphrase.

Once you entered your passphrase twice, click on \Button{Next}.
%%Oder lieber so:  Dann mussen wir den Button in GPA aber �ndern.
%%Fertig? Dann klicken Sie auf \Button{Fertig} und Ihre Schl�ssel werden
%%endg�ltig erzeugt.


\clearpage
%% Original page 21

This starts the creation of your key pair, which can take a few minutes.
In the meantime you can continue to work with other programs on your
computer, which will slightly increase the quality of the key being
generated.

%put a note in here somewhere that 'key pair' and 'key' are used
interchangeably.

Once the key is generated, the following window appears:

% screenshot: Wizard, backup
\begin{center}
\IncludeImage[width=0.6\textwidth]{sc-gpa-gen-backup}
\end{center}

This window asks you to create a back-up copy of your key. Please do
this now, even if this is just a test run.

If the default file name is acceptable, click on  \Button{OK}. If you
would like to store the back-up in a different location, please select a
different file name now. 

\textbf{Important:} Once the the back-up file is saved onto your hard
drive, you should save this file onto a removable file storage device
(ex. USB stick, diskette or CD-ROM), and subsequently delete the
original file.  Be sure to keep the file storage device in a secure
place. 

Please note that you can create a back-up copy at any time, by selecting
the following option from the main menu: 
\Menu{Key$\rightarrow$Backup}.


\clearpage
%% Original page 23
\textbf{
This concludes the installation of Gpg4win and the generation of your
key pair.  You are now the owner of a unique and secure digital key. 
}

You should now see the main page of GPA. The key pair you just created
will be shown in the middle of the window, beside the key pair symbol.  

Clicking on the key pair symbol will provide details about your key
pair, which are examined in more detail below.

% screenshot: keymanger mit einem Schl�ssel.
\begin{center}
\IncludeImage[width=0.9\textwidth]{sc-gpa-first-key}
\end{center}

What is the significance of the notes regarding your key? Your key has
no built-in expiry date, and is therefore valid indefinitely.  You can
however change the validity period of the key - more about this later.

A key consisting of 1024 bits is considered very secure, without
creating undue strain on your computer's capacity.


\textbf{ \OtherBook{} More information on this topic can be found in
Chapter 5 "`More about keys"' in the manual "`\xlink{Gpg4win for
Advanced Users}{durchblicker.html}"'.}

\clearpage
%% Original page 24

\section{Publishing your key per  \Email{}}

One of the more practical aspects of Gpg4win is its ability to use a
"`non-secret"' public key for the en- and decryption of data.  As long
as your key and its corresponding passphrase are secure, you have gone a
long way towards keeping your information confidential. 

Everyone can and should have your public key, just as you can and should
have the public keys from those you correspond with. 

Because:

\textbf{
In order to exchange secure \Email{}s, each party must have access to
the public key of the other party. 
}

So in order to send someone encrypted \Email{}s, you must have their
public key in order to encrypt the emails sent to them. 

Similarly, someone wanting to send you encrypted \Email{}s must also
have your public key in order to encrypt the email being sent to you.  

It is for this reason that your public key should be made widely
available. Depending on the number of people on your mailing list, this
can be done in two ways:

\begin{itemize}
\item \textbf{directly sending an \Email{}} to selected recipients
\item \textbf{or publishing the key on a key server} --- making it
available to anyone
\end{itemize}

The first way of distributing your public key is to send it per
\Email{} to one or more selected recipients.
Alternatively you can make your \Email{}address available to anyone
throughout the Internet. The second option is somewhat risky, as it can
result in considerable SPAM activity on your email account. Therefore it
is a good idea to only use an address with a good SPAM filter. 


\clearpage
%% Original page 25
You can use Adele to practice the following steps:

Adele is a very good \Email{}robot for practicing secure correspondence.
And because most people prefer to correspond with a real person rather
than with a piece of software (which is what Adele is, after all), we
developed the following scenario:

% Cartoon:  Adele mit Buch ind er Hand vor Rechner ``you have mail"'
\begin{center}
\IncludeImage{adele01}
\end{center}

You first send Adele your public key. Once Adele has received your key,
she uses it to encrypt an  \Email{} which she sends to you along with
her own public key.  You can now decrypt Adele's email using your own
private key, and you can also respond to Adele by encrypting your email
with her public key.

You can now export your public key, copy it into an \Email{} und send it
to Adele.


\clearpage
%% Original page 26

Here is one possible way of doing just that, a method which works even
if your \Email{} service does not allow attachments.
This procedure will also provide you with a first in-depth look at your
key and its components.

How it works:\\
Select the key you want to export by clicking on the corresponding key
on your list, and then clicking on the \Button{Export} icon of the main
GPA menu. Choose a file to export your key to, ex. \\
\Filename{my-key.asc}.  A popup window will let you know whether the
operation was successful.  Then click on \Button{OK}.


\clearpage
%% Original page 27

You can access the file thru Windows Explorer, making sure you choose
the same folder you chose when exporting the key to file.  You can open
the file with a text editor (ex. WordPad), which will show your public
key as a series of blocks containing text and numbers. 

% screenshot: Editor mit ascii armored key
\begin{center}
\IncludeImage[width=0.9\textwidth]{sc-misc-mein-key-asc}
\end{center}


\clearpage
%% Original page 28

Highlight the whole key portion from 

\begin{verbatim}
-----BEGIN PGP PUBLIC KEY BLOCK-----
\end{verbatim}
to
\begin{verbatim}
-----END PGP PUBLIC KEY BLOCK-----
\end{verbatim}

and copy it using the copy function on your toolbar or a keyboard
shortcut such as Ctrl-C. This saves your key on the clipboard until you
are ready to paste it, as described below. 

Start your email program. Then open a new \Email{} message and paste
your public key (Windows users may use a shortcut key such as Ctrl-V).
Prior to doing this, it is useful to configure your email program to
send messages in text-only format, rather than HTML.

Put \verb-adele@gnupp.de- into the address line of the \Email{}, and
\textbf{my public key} on the subject line.

Your \Email{} should look like this:

% screenshot:  Eines composer Windows.
\begin{center}
\IncludeImage[width=0.9\textwidth]{sc-ol-send-test-key}
\end{center}

Now send this \Email{} to Adele. Make sure that you include your proper
email address as sender rather than your practice address, otherwise you
will never get a response from Adele!

\clearpage
%% Original page 29

This process works exactly the same way if you send your key to a real
\Email{} address. You can also add comments, just like any other emails.
Obviously this is not required for emails addressed to Adele, as the
robot's only purpose is to assist you with the technical aspects of this
process.

\textbf{Summary:} You have now sent your public key per \Email{} to
someone else (ex. Adele). 

\textbf{ \OtherBook{} The "`Copy \& Paste"'-method demonstrated in this
example is a more transparent and easier-to-understand method for
beginners. Chapter 7 of the manual "`\xlink{Gpg4win for Advanced
Users}{durchblicker.html}"' describes how you can send your key as a
file attachment, which is a more commonly-used method.}  

\clearpage
%% Original page 30
\section{Sending your key to a keyserver}

You can use this option even if you are exchanging encrypted \Email{}s
with only a few people. By making your key available on a keyserver, it
is always accessible to you and/or others.  

\textsc{PLEASE NOTE: Although there are no indications that spammers are
currently collecting addresses from keyservers, it is technically
possible to do so. Therefore we do not recommend publishing your key to
a key server if you do not have an effective spam filter.}

Choose your key by clicking on it. Then click on \Menu{Server
$\rightarrow$Send key} which is found under \Menu{Server$\rightarrow
$Server}.  A popup window will alert you that your key is about to be
sent to a public key server. 

% screenshot: GPA export key to keyserver
\begin{center}
\IncludeImage[width=0.4\textwidth]{sc-gpa-ks-export-p}
\end{center}

This window also displays a default key server.  If you click on
\Button{Yes}, your key will be automatically sent to the server, from
where it is sent to a network of worldwide key servers. Now anyone can
download your key, and use it to send you a secure \Email{}.

If this is a practice run, do NOT send your practice key to the key
server, because it cannot be removed once sent.

\textbf{Summary:} You now know how to distribute your key via a key
server on the Internet. 

\textbf{ \OtherBook{} Chapter 6 of the manual"`\xlink{Gpg4win for
Advanced Users}{durchblicker.html}"'explains how to find someone else's
key on a keyserver.}

\clearpage
%% Original page 31
\section{Decrypting an \Email{}}

Adele has received your public key, and uses it to encrypt an \Email{}
which she sends back to you. 

% cartoon: Adele typing and sending a mail
\begin{center}
\IncludeImage{adele02}
\end{center}


\clearpage
%% Orginal page 32
The email will look like this:

\begin{verbatim}
From: Adele (The friendly email-Robot) <adele@gnupp.de>
Subject: Re: my publis key
To: larrys@gpg4win.de
Date: Thu, 12 Jan 2006 09:17:28 +0100

-----BEGIN PGP MESSAGE-----
Version: GnuPG v1.4.1 (GNU/Linux)

hQEOA9FS8I3hSvdPEAP/W6W6f4MBwqTdzd9O/7FOTDHh//bQ+GUWoT0k9Y0i96UZ
QO1VhQSia6a8DZrFQj7SlJWmB1MM7RNhkmhfZsD5Bn9ICmwwOt2xJDBkCQ34gu5N
NxQ92WXZjHCaI0dSlynNziNbK8Ik26YPBYkQjLUDhHN4CRZ7q67eVEd/B9DI04wD
....
ujbjyj09L/9NvoBniWrgqVUayKr1Ls8OIZkyiex6mKypPGADJFAzvTwjubj5S6zJ
A+QvSXUB9Hj8Ft2Nt3j0B/gWn5no3Er2/15UcBn/UPSxW9or0w9seDxCuSXvpakX
bcneOm/pcJNEHcApXWXpoNOxRZ1MksM300w+79M6p2w=
=VCHb
-----END PGP MESSAGE-----
\end{verbatim}

(Please note that this encryption sequence has been shortened
considerably for the purposes of this manual.)
\clearpage
%% Original page 33

\textbf{Using WinPT to decrypt this \Email{}.}

WinPT (Windows Privacy Tray key management tool) serves as PnuPG's
"`front end"'. This program is used to en- and decrypt \Email{}s; it
also creates and verifies digital signatures.  One of WinPT's advantages
is that it works with any email program. 

Most email programs (ex. MS Outlook for Windows) also have special
plug-ins which can en- and decrypt email directly within the program.


\textbf{ \OtherBook{} Chapter 8 of the manual "`\xlink{Gpg4win for
Advanced Users}{durchblicker.html}"' provides more information on this
particular funtion.} 

WinPT performs the en/decryption function using your computer's memory.
This means that any text to be en/decrypted should be copied onto the
computer's clipboard.

To do this, highlight the complete text in Adele's \Email{} and copy it
onto the Clipboard using the copy function or the shortcut Ctrl-C. 


\clearpage
%% Original page 34
Starting WinPT from Windows Start Menu:

% screenshot: startmenu Auswahl winpt
\begin{center}
\IncludeImage[width=0.7\textwidth]{sc-winpt-startmenu}
\end{center}

% the english version of the program does not have the "hinweis"below,
so I have not translated it.

W�hrend WinPT startet, erscheint kurz ein Hinweis darauf, dass das
Programm die bereits vorhandenen Schl�ssel einl�dt.

Once you open the program, the WinPT icon - a key - appears in your
Windows taskbar located in the lower right hand corner of the screen.

% screenshot: taskbar rechts mit WinPT icon
\begin{center}
\IncludeImage{sc-winpt-trayicon}
\end{center}


\clearpage
%% Original page 35
Right-click on the icon to open the menu box, and click on
\Menu{Clipboard$\rightarrow$Decrypt/Verify}.

The program will then ask you for your secret passphrase and proceed to
decrypt Adele's \Email{}.

% screenshot: WinPT decryption dialog
\begin{center}
\IncludeImage[width=0.7\textwidth]{sc-winpt-clip-decrypt}
\end{center}

A popup window window will notify you when the decryption process is
complete. 

\clearpage
%% Original page 36

The decrypted text is now located on the clipboard, much in the same way
as was done during the encryption process. You can copy the text (using
shortcut key Ctrl-V) into your text editor or email program. 

Adele's decrypted email will look something like this\footnote{Depending
on Adele's software version, it might look slightly different.}:

\begin{verbatim}
Hello Larry Smith,

here is an encrypted response to your email. 

I received the your public key with the ID 57251332CD8687F6 and a
description `<larrys@gpg4win.de>'.

I have enclosed the public key of adele@gnupp.de.

Sincerely,
adele@gnupp.de
\end{verbatim}

The text sequence following this message is Adele's public key. 

Below we show you how to import someone's public key and attach it to
your key ring, which allows you to use this key anytime to decrypt
messages from that person (in this case, Adele). 

\textbf{Summary:}

\begin{enumerate}
\item You know how to decrypt an encrypted \Email{} using your private
key.

\item You know how to send an encrypted response using the public key of
the sender. 
\end{enumerate}


\clearpage
%% Original page 37
\section{Attaching a key to your key ring}

You can store your friends' public keys on your "`key ring"' so they do
not have to send you their public key every time they write to you.  

\textbf{First Option:}

To import a public key (ie. attach it to your key ring), you can save
the key as a text block, much like you did with your own key. 


To do this:

Highlight the public key portion of the email you received, from

\begin{verbatim}
-----BEGIN PGP PUBLIC KEY BLOCK-----
\end{verbatim}
to
\begin{verbatim}
-----END PGP PUBLIC KEY BLOCK-----
\end{verbatim}

and use Copy \& Paste to insert it into your text editor.  We recommend
using a file and folder name that is easily found later; ex.
\Filename{adeles-key.asc} in folder 
\Filename{My Documents}.



\clearpage
%% Original page 38
\textbf{2nd Option:}

The key is sent to you as a file attachment to the \Email{}. No matter
which mail program you use, you can always save attachments onto your
hard drive.  Do this now (again using names you will easily recognize
and find later on, ex. \Filename{My Documents}).

It does not matter whether you save the key as text or directly as an
\Email{}attachment, as each method imports the key into your GnuPG-"`Key
Ring"'.

% screenshot:  Startmenu, Auswahl GPA

This is how it works:

Start the GNU Privacy Assistant (GPA) from Windows (this is necessary
only if you shut it down during the previous practice session).


\clearpage
%% Original page 39
Click on Import, then select and load the key file.

% screenshot GPA Schl�sselverwaltung mit zwei Schl�sseln
\begin{center}
\IncludeImage[width=0.9\textwidth]{sc-gpa-two-keys}
\end{center}

You have now imported someone else's (in this case Adele's) public key
and attached it to your key ring. Now you can use this key to send
encrypted messages to the owner of that public key, as well as to verify
his or her signature.


\textbf{Before continuing, it is important to address the following
concern:}\\
It is possible that the \Email{} was sent by someone else using Adele's
name, therefore how do you know that the public key sent to you is
really Adele's key? 

\textbf{ \OtherBook{} Chapter 9 ("`Key Verification"') of the manual "`
\xlink{Gpg4win for Advanced Users}{durchblicker.html}"' addresses this
important question.  You may want to read that section now before
continuing with this manual. }

\clearpage
%% Original page 40
Chapter 9 of the manual "`Gpg4win for Advanced Users"' shows you how to
validate a key as well as to sign a message (ie. attach a signature)
using your private key.

Chapter 10 of the advanced manual also discusses ways to attach a
signature to \Email{}messages. This is the equivalent of attaching an
electronic seal to your message, which allows the recipient to verify
whether the \Email{} has been altered during transmission. 

The signature verification process is fairly simple. To do this, you
must have the sender's public key on your Gpg4win-"`key ring"' (see
Chapter 8 of "`Gpg4win for Advanced Users"' for more information).

% cartoon:  M�ller mit Schl�ssel
\begin{center}
\IncludeImage[width=0.4\textwidth]{man-with-signed-key}
\end{center}


\clearpage
%% Original page 41
You can tell whether an \Email{} has been electronically signed if the
text of the email is framed with the sender's signature (like a border).
It will look something like this: 

\begin{verbatim}
-----BEGIN PGP SIGNED MESSAGE-----
Hash: SHA1
\end{verbatim} 

and ends \Email{}-message with 

\begin{verbatim}
-----BEGIN PGP SIGNATURE-----
Version: GnuPG v1.4.2 (MingW32)

iEYEARECAAYFAjxeqy0ACgkQcwePex+3Ivs79wCfW8u
ytRsEXgzCrfPnjGrDDtb7QZIAn17B8l8gFQ3WIUUDCMfA5cQajHcm
=O6lY
-----END PGP SIGNATURE-----
\end{verbatim}

Highlight the text starting from \textsl{BEGIN PGP SIGNED MESSAGE} to
\textsl{END PGP SIGNATURE} and copy it (using Ctrl-C) to your
clipboard. 

Now cotinue to decrypt the \Email{} as shown in Chapter 7 of this
manual.

Right-click on the WinPT icon on your Windows taskbar and select
\newline
\Menu{Clipboard$\rightarrow$ Decrypt/Verify}.

You should see the following window:

% screenshot: WinPT good signature.
\begin{center}
\IncludeImage[width=0.7\textwidth]{sc-winpt-good-sig}
\end{center}

If the status line on the window displays the following message
\emph{Invalid Signature}%someone check if this is the right text used in
the program, i did not encounter it so I am not sure of the tranlsation,
although i think this should be right

it means that the message has been altered after it was sent.  This does
not always mean that a third party has altered the message, as it could
also have been altered through a technical error during transmissiont
through the Internet. 

\textbf{ \OtherBook{} Before continuing, you may want to read Chapter 10
of the manual "`\xlink{Gpg4win for Advanced Users}{durchblicker.html}"'
which contains additional information on how to handle and manage
invalid signatures.}

\clearpage
%% Original page 42
\section{Encrypting \Email{}s}

Encrypt an \Email{} and send it to Adele (or another person).

Open your \Email{}program and write a message (if you are sending the
message to Adele, any text will do)$\ldots$

Highlight the text and copy it onto your Clipboard (Windows) using the
copy function or the Ctrl-C shortcut.

% screenshot: composer mit mail an Adele, body markiert
\begin{center}
\IncludeImage[width=0.9\textwidth]{sc-ol-send-enc-msg1}
\end{center}

Right click on the WinPT icon on your Windows Taskbar and select
\Menu{Clipboard$\rightarrow$Encrypt}.


\clearpage
%% Original page 45
You should now see a window containing the keys on your key ring. Using
the examples featured in this manual, you would see Adele's key (the one
she sent to you), as well as your own key, which you created in Chapter
2.

% screenshot: Winpt encryption key selection
\begin{center}
\IncludeImage[width=0.7\textwidth]{sc-winpt-sel-enc-key}
\end{center}

Click on Adele's key to encrypt the message you just composed, keeing in
mind the following principle:

\textbf{In order to send someone encrypted \Email{}s, you need to have
that person's public key to encrypt the message properly.}

Clicking on \Button{OK} starts the encryption process, which will be
confirmed by the program if successful.

\clearpage
%% Original page 46
The encrypted message can be found on your computer's clipboard, from
which you can copy it into the \Email{} window.  You can erase the
unencrypted text or just copy and paste the contents of the clipboard.

It would look something like this: 

% screenshot: composer Windows mit enrcypted text 
\begin{center}
\IncludeImage[width=0.9\textwidth]{sc-ol-send-enc-msg2}
\end{center}

Now send your \Email{} to Adele.  Remember to use your proper email
address (rather than the previous practice email) otherwise you will not
receive a response.   $\ldots$

\textbf{
Congratulations! You have just encrypted your first \Email{}!
}

\clearpage
%% Original page 47
\section{How to archive/store encrypted \Email{}s}

It is probably not wise to store the full(ie. decrypted) text of your
encrypted \Email{}s on your computer (they were encrypted for a reason,
after all). Therefore it is advisable to always store the encrypted
copies of the \Email{}.

You may already guess the problem at hand: In order to decrypt the
archived \Email{}s you need the private key of the recipient, which will
(or at least should) never be available to you. $\ldots$

The solution: You encrypt messages to yourself as well. 

You can encrypt a message designated for a recipient (ex. Adele) using
their own as well as your own public key. This allows you to decrypt the
text later using your private key.

Since Gpg4win cannot tell the difference between the different keys used
to encrypt a message (ie. you could have more than one key), you need to
set up the program accordingly. 

\clearpage
%% Original page 48
To use this option, right click on the WinPT icon select
\Menu{Preferences $\rightarrow$ GPG}.

% screenshot: Winpt configuration dialog
\begin{center}
\IncludeImage[width=0.7\textwidth]{sc-winpt-enctoself}
\end{center}

This opens the GnuPG Preferences Window. You can add your key in the
"Encrypt to this key" field using your \Email{}address.

\Email{}programs which directly support GnuPG will also feature this
option.

\clearpage
%% Original page 49
\textbf{Summary:}

\begin{enumerate}
\item You responded to the recipient by encrypting an \Email{} with
his/her public key.

\item You also set up WinPT to encrypt archived copies of your \Email{}s
with your own private key.
\end{enumerate}

\textbf{And that's it - welcome to the world of free and secure \Email{}
encryption!}

\textbf{
\OtherBook{} We recommend you read Chapter 10 to 12 of the manual "`
\xlink{Gpg4win
for Advanced Users}{durchblicker.html}"' for additional information on
adding signatures to \Email{}s, as well as importing and using an
existing private key in GnuPG.
}

\textbf{
\OtherBook{} Furthermore, Chapters 13 and 14 of the manual "`
\xlink{Gpg4win for Advanced Users}{durchblicker.html}"' provide you with
additional information regarding the premises on which GnuPG's security
features are based, as well as detailed insights into the mathematical
concepts behind GnuPG.  These texts are easily readable, as they were
created for \textbf{the every day user}, not mathematicians or
cryptographs. }


\newpage
\appendix
\section{Suggestions regarding the Outlook plugin \emph{GPGol}}

GPGol, a Microsoft Outlook plugin, integrates the operation of GnuPG.
Here are some tips regarding the operation of this plugin.

While it is fairly easy to operate the plugin in other \Email{}programs,
the technical difficulties associated with the integration of OpenPGP in
Outlook make it harder to operate the plugin in this context.

The current version of GPGol works only with Outlook 2003 SP2, and will
notify you if you are using an older version of Outlook. 

You should also be aware of the following limitations:

\begin{itemize}
\item Do not use Word as your text editor. 

\item Inline-PGP or the conventional PGP is fully supported, however,
PGP/MIME \Email{}s cannot be created. 

On the other hand, it does support the decryption and signature
verification of simple (unlike more complex) PGP/MIME \Email{}s.
\end{itemize}


\subsection{Installation}

The plugin can be installed using the the Gpg4win Installer. The next
time you start Outlook, a \Menu{GnuPG} tab will be added in the
\Menu{Tools$\rightarrow$Options} menu :

\begin{center}
\IncludeImage[width=0.7\textwidth]{sc-gpgol-options}
\end{center}

The first two options of this tab allow you to set default settings
regarding the encryption and signature of new messages. Note that you
can always change these settings at the time you are creating new
messages.

The option "Save decrypted message automatically" will save attachments
as unencrypted text (after decryption). As a result, the attachment can
be read anytime without having to decrypt it again. 

The option "Automatically sign attachments" ensures that attachments are
automatically signed along with your main text. For this purpose, an
attachment with a signature is created for every attachment included in
the message.

The option "Also encrypt message with the default key" allows you to
enter the Key ID of your own key so that your messages are automatically
(and additionally) encrypted using your own key. This allows you to
later retrieve, decrypt and read messages in the Sent folder.

The option "`Also decrypt in preview window"' is only feasible with very
fast computers; at the present time it has only limited operational
capacity. 

The option "`Show HTML view if possible"' enables you to view a message
in HTML format. In most cases, or when an HTML format is not available,
it will be shown in text format. 

There are a number of preset default settings which take effect after
the initial installation. However, you should also ensure that you are
not using \textbf{Microsoft Word} to compose messages.  We also
recommend against using HTML messages.

You can set these option using the \Menu{Mail Format}tab in the
\Menu{Tools$\rightarrow$Options} menu.  Your settings should look like
this:

\begin{center}
\IncludeImage[width=0.7\textwidth]{sc-gpgol-noword}
\end{center}


\subsection{Common Questions}

\begin{description}
\item[Encryption Icons are not visible in the notification area/system
tray.] 

%The german text below was somewhat confusing - I know what they are
trying to accomplish here and changed the english text slightly.

If the system tray contains a large number of icons, it may not display
all of them directly. Clicking on the arrow at the end of the system
tray (located on the bottom right hand side of most computers) will open
the tray to show all icons including the Key Manager icon.

Again, make sure that you are not using Microsoft Word to compose
messages (see installation instructions).
 
\item[What is the meaning of the letters included in the GnuPG tab?] 

  GPGol lists the components and their associated file names as part of
the decryption or signature verification of messages. File names
containing encrypted attachments are prefaced with "`E"', while signed
attachments are marked with an "`S"'. 

\item[Where can I locate information regarding the latest version
GPGol?]       

  Click on the logo on the lower left side of your GnuPG options menu.

\item[Why is it not possible to cancel an encryption process?] 

  GPGol begins the encryption process as soon as you press the Send
button. An error contained in Outlook prevents you from interrupting or
cancelling this process. To avoid this, we recommend configuring Outlook
so that messages are not sent immediately. This gives you the
opportunity to cancel the message before it is sent. As a security
measure, GPGol will try to erase the contents of the message if the
encryption process is interrupted, but this is not always successful.

\item[Why does the confirmation window appear when accessing certain
\Email{}s?]

If \linebreak{} GPGol is not installed as a trusted plugin, Outlook
assumes that it is trying to access internal information without
authorization.  GPGol attempts to prevent this from happening, but
sometimes it is necessary in order to properly display encrypted or
signed \Email{}s.

GPGol is still in the development process. One outstanding issue is its
registration as a trusted plugin. Future versions ofGPGol should be able
to solve the issue of having to confirm each or certain emails.

\item[How come GPGol cannot create PGP/MIME messages?]

At the present there is no way of telling Outlook that a PGP/MIME
message is to be created. Outlook decides on its own which
"`Content-Type"'is to be used, as the plugin is not able to preset a
certain type.  You may want to contact Microsoft directly to talk about
this setting and/or missing documentation regarding this setup.

\item[Why are signature verifications not performed automatically?]

We are currently working on an automatic signature verification option
to be used when opening an \Email{}.  However, given Outlook's complex
technical environment, this is not an easy issue to solve. 


\end{description}

\clearpage
\section{Transferring from other GnuPG programs}
\label{ch:migration}

This section explains how to transfer to Gpg4win from other GnuPG-based
programs. The installation program will recognize some of these programs
and notify you if that is the case.

As a general rule it is recommended that an existing GnuPG-based program
be removed before installing Gpg4win. 

Do not forget to save existing keys. The best way to do this is to use
the options in the existing system. Select the option to save a private
(secret) key, as well as all existing public keys.  Save these in one or
two files.

As soon as Gpg4win is installed, check if your existing keys are listed
in the new program (using either GPA or WinPT). If yes, the previous
system was able to use the new rules regarding key storage, and no
additional action is required on your part. 

If the program does not list the existing keys, you can import them from
your backup file.  For more information on this topic, refer to Chapter
12 in the manual "`\xlink{Gpg4win for Avanced
Users}{durchblicker.html}"'.

If your older system also includes GPA, you can use its back-up option
which should be very similar to the operation of GPA in Gpg4win.

Also, if you are not able to locate your existing keys, use the 'Search'
function in Windows to find the files named \Filename{secring.gpg} and
\Filename{pubring.gpg} an import them per GPA.\footnote{This is not the
official procedure, but it works with all current versions of GnuPG.}




\clearpage
\section{History}

\begin{itemize}
\item  "`GnuPP for Beginners"', First Edition March 2002,\\
  Authors: Manfred J. Heinze, TextLab text+media\\
  Consulting: Lutz Zolondz, G-N-U GmbH\\
  Illustrations: Karl Bihlmeier, Bihlmeier \& Kramer GbR\\
  Layout: Isabel Kramer, Bihlmeier \& Kramer GbR\\
  Documentation: Dr. Francis Wray, e-mediate Ltd.\\
  Editor: Ute Bahn, TextLab text+media\\
  Published by the Federal Ministry for Industry and Technology.\\
  Available at
  \verb-http://www.gnupp.de/pdf/einsteiger.pdf-.
%    Der Abschnitt "`History"' ist im Originaldokument nicht vorhanden
%    und wurde von Werner Koch beigef�gt.
\item Revised unpublished version by TextLab text+media.      
\item "`Gpg4win for Beginners"', December 2005\\
      Authors:  Werner Koch, g10 Code GmbH\\
      Published as part of the Gpg4win project.
\item "`Gpg4win for Novices"', November 2006\\
      Translated by: Brigitte Hamilton\\
      Published as part of the Gpg4win project.

\end{itemize}

\clearpage
\input{fdl.tex}

\end{document}
